\documentclass[11pt,a4paper]{article}
\usepackage[utf8]{inputenc}
\usepackage[margin=2cm]{geometry}
\usepackage{amsmath,amssymb}
\usepackage{enumitem}
\usepackage{fancyhdr}
\usepackage{tabularx}
\usepackage{booktabs}
\usepackage{xcolor}

\setlength{\headheight}{14pt}
\pagestyle{fancy}
\fancyhf{}
\lhead{\textbf{MA1101 Matematika I --- Solusi \& Rubrik Kuis (20 Menit)}}
\rhead{\textbf{Semester 1 TA 2026/2027}}
\rfoot{Halaman \thepage}

\definecolor{ITBblue}{RGB}{0, 75, 141}

\begin{document}

\begin{center}
    {\large \textbf{FAKULTAS MATEMATIKA DAN ILMU PENGETAHUAN ALAM}}\\
    {\large \textbf{INSTITUT TEKNOLOGI BANDUNG}}\\[4pt]
    {\large \textbf{\color{ITBblue}KUNCI JAWABAN \& PEDOMAN PENSKORAN KUIS KELAS}}\\[2pt]
    \textbf{Materi: Subbab 2.1--2.7, 2.9, 3.1--3.3 \quad | \quad Total Skor: 100 Poin}
\end{center}

\vspace{6pt}
\hrule
\vspace{10pt}

% =========================================================================
% SOLUSI PAKET A
% =========================================================================
\section*{Solusi \& Rubrik Paket A}

\subsection*{Soal 1 [Bobot: 45 Poin]}

\paragraph{(a) [20 Poin] Nilai $f'(1)$ dari $f(x) = x^2 + 2x$ Menggunakan Definisi Limit}
\begin{itemize}[leftmargin=18pt]
    \item \textbf{Langkah 1: Menuliskan rumus definisi limit turunan di $a = 1$}
    \[
    f'(1) = \lim_{h \to 0} \frac{f(1+h) - f(1)}{h} \quad \text{\textbf{[4 poin]}}
    \]
    \item \textbf{Langkah 2: Menghitung $f(1)$ dan $f(1+h)$}
    \begin{align*}
        f(1) &= 1^2 + 2(1) = 3 \\
        f(1+h) &= (1+h)^2 + 2(1+h) = 1 + 2h + h^2 + 2 + 2h = h^2 + 4h + 3 \quad \text{\textbf{[5 poin]}}
    \end{align*}
    \item \textbf{Langkah 3: Menghitung selisih pembilang}
    \[
    f(1+h) - f(1) = (h^2 + 4h + 3) - 3 = h^2 + 4h = h(h + 4) \quad \text{\textbf{[4 poin]}}
    \]
    \item \textbf{Langkah 4: Membagi dengan $h$ ($h \neq 0$)}
    \[
    \frac{f(1+h) - f(1)}{h} = \frac{h(h+4)}{h} = h + 4 \quad \text{\textbf{[4 poin]}}
    \]
    \item \textbf{Langkah 5: Evaluasi limit saat $h \to 0$}
    \[
    f'(1) = \lim_{h \to 0} (h + 4) = 4 \quad \text{\textbf{[3 poin]}}
    \]
\end{itemize}

\paragraph{(b) [25 Poin] Turunan Implisit, Garis Singgung, dan Hampiran Linear}
Kurva: $x^2 + \sin(y) = 4$ di titik $(2, 0)$.
\begin{enumerate}[label=(\roman*), leftmargin=18pt]
    \item \textbf{[10 Poin] Menentukan $\dfrac{dy}{dx}$ di $(2, 0)$ secara implisit:}
    \begin{align*}
        \frac{d}{dx}\left[ x^2 + \sin(y) \right] &= \frac{d}{dx}[4] \\
        2x + \cos(y) \frac{dy}{dx} &= 0 \quad \text{\textbf{[5 poin: turunan $x^2$ dan aturan rantai $\sin(y)$]}} \\
        \frac{dy}{dx} &= \frac{-2x}{\cos(y)} \quad \text{\textbf{[2 poin]}}
    \end{align*}
    Substitusikan titik $(2, 0)$:
    \[
    \left.\frac{dy}{dx}\right|_{(2, 0)} = \frac{-2(2)}{\cos(0)} = \frac{-4}{1} = -4 \quad \text{\textbf{[3 poin]}}
    \]

    \item \textbf{[7 Poin] Persamaan garis singgung di $(2, 0)$:}
    Gradien $m = -4$.
    \begin{align*}
        y - y_0 &= m (x - x_0) \\
        y - 0 &= -4 (x - 2) \implies y = -4x + 8 \quad \text{\textbf{[7 poin]}}
    \end{align*}

    \item \textbf{[8 Poin] Taksiran nilai $y$ ketika $x = 2{,}01$:}
    \begin{itemize}
        \item \textit{Melalui Hampiran Linear (Garis Singgung):}
        \[
        L(x) = -4x + 8 \implies y(2{,}01) \approx L(2{,}01) = -4(2{,}01) + 8 = -8{,}04 + 8 = -0{,}04 \quad \text{\textbf{[8 poin]}}
        \]
        \item \textit{Melalui Diferensial:}
        $\Delta x = dx = 2{,}01 - 2 = 0{,}01$.
        \[
        dy = \left(\left.\frac{dy}{dx}\right|_{(2, 0)}\right) dx = (-4)(0{,}01) = -0{,}04
        \]
        Maka $y(2{,}01) \approx y_0 + dy = 0 + (-0{,}04) = -0{,}04$.
    \end{itemize}
\end{enumerate}

\vspace{10pt}
\subsection*{Soal 2 [Bobot: 55 Poin]}
Fungsi: $f(x) = x^3 - 3x$ pada selang tertutup $[-2, 2]$.

\paragraph{(a) [20 Poin] Kemonotonan \& Ekstrem Lokal}
\begin{itemize}[leftmargin=18pt]
    \item \textbf{Turunan pertama:}
    \[
    f'(x) = 3x^2 - 3 = 3(x^2 - 1) = 3(x-1)(x+1) \quad \text{\textbf{[5 poin]}}
    \]
    Titik stasioner: $x = -1$ dan $x = 1$.
    \item \textbf{Uji tanda $f'(x)$:}
    \begin{itemize}
        \item Pada $(-2, -1)$: $f'(x) > 0 \implies f$ \textbf{monoton naik} pada $[-2, -1]$.
        \item Pada $(-1, 1)$: $f'(x) < 0 \implies f$ \textbf{monoton turun} pada $[-1, 1]$.
        \item Pada $(1, 2)$: $f'(x) > 0 \implies f$ \textbf{monoton naik} pada $[1, 2]$.
    \end{itemize}
    (\textit{Menyatakan selang kemonotonan dengan benar: \textbf{[8 poin]}})
    \item \textbf{Ekstrem lokal:}
    \begin{itemize}
        \item Di $x = -1$: tanda $f'$ berubah dari $+$ ke $-$, maka ada \textbf{maksimum lokal} di $x = -1$ dengan nilai $f(-1) = (-1)^3 - 3(-1) = 2$. \textbf{[3.5 poin]}
        \item Di $x = 1$: tanda $f'$ berubah dari $-$ ke $+$, maka ada \textbf{minimum lokal} di $x = 1$ dengan nilai $f(1) = 1^3 - 3(1) = -2$. \textbf{[3.5 poin]}
    \end{itemize}
\end{itemize}

\paragraph{(b) [15 Poin] Kecekungan \& Titik Belok}
\begin{itemize}[leftmargin=18pt]
    \item \textbf{Turunan kedua:}
    \[
    f''(x) = \frac{d}{dx}[3x^2 - 3] = 6x \quad \text{\textbf{[4 poin]}}
    \]
    $f''(x) = 0 \iff x = 0$.
    \item \textbf{Selang Kecekungan:}
    \begin{itemize}
        \item Pada $(-2, 0)$: $f''(x) < 0 \implies$ kurva \textbf{cekung ke bawah}. \textbf{[3 poin]}
        \item Pada $(0, 2)$: $f''(x) > 0 \implies$ kurva \textbf{cekung ke atas}. \textbf{[3 poin]}
    \end{itemize}
    \item \textbf{Titik Belok:}
    Karena kecekungan berubah tanda di $x = 0$, terdapat titik belok di $(0, f(0)) = (0, 0)$. \textbf{[5 poin]}
\end{itemize}

\paragraph{(c) [20 Poin] Titik Kritis, Nilai Maksimum Global, \& Nilai Minimum Global pada $[-2, 2]$}
\begin{itemize}[leftmargin=18pt]
    \item \textbf{Titik-titik kritis pada $[-2, 2]$:}
    \begin{itemize}
        \item Titik ujung selang: $x = -2, 2$.
        \item Titik stasioner: $x = -1, 1$.
        \item Titik singular: tidak ada.
    \end{itemize}
    Daftar titik kritis: $\{-2, -1, 1, 2\}$. \textbf{[6 poin]}
    \item \textbf{Evaluasi nilai fungsi pada seluruh titik kritis:}
    \begin{align*}
        f(-2) &= (-2)^3 - 3(-2) = -8 + 6 = -2 \quad \text{\textbf{[2 poin]}} \\
        f(-1) &= 2 \quad \text{\textbf{[2 poin]}} \\
        f(1)  &= -2 \quad \text{\textbf{[2 poin]}} \\
        f(2)  &= 2^3 - 3(2) = 8 - 6 = 2 \quad \text{\textbf{[2 poin]}}
    \end{align*}
    \item \textbf{Kesimpulan Ekstrem Global:}
    \begin{itemize}
        \item \textbf{Nilai Maksimum Global} adalah \textbf{2} (tercapai di $x = -1$ dan $x = 2$). \textbf{[3 poin]}
        \item \textbf{Nilai Minimum Global} adalah \textbf{-2} (tercapai di $x = -2$ dan $x = 1$). \textbf{[3 poin]}
    \end{itemize}
\end{itemize}

\newpage

% =========================================================================
% SOLUSI PAKET B
% =========================================================================
\section*{Solusi \& Rubrik Paket B}

\subsection*{Soal 1 [Bobot: 45 Poin]}

\paragraph{(a) [20 Poin] Nilai $g'(1)$ dari $g(x) = x^2 + 3x$ Menggunakan Definisi Limit}
\begin{itemize}[leftmargin=18pt]
    \item \textbf{Langkah 1: Rumus definisi limit:}
    \[
    g'(1) = \lim_{h \to 0} \frac{g(1+h) - g(1)}{h} \quad \text{\textbf{[4 poin]}}
    \]
    \item \textbf{Langkah 2: Menghitung $g(1)$ dan $g(1+h)$:}
    \begin{align*}
        g(1) &= 1^2 + 3(1) = 4 \\
        g(1+h) &= (1+h)^2 + 3(1+h) = 1 + 2h + h^2 + 3 + 3h = h^2 + 5h + 4 \quad \text{\textbf{[5 poin]}}
    \end{align*}
    \item \textbf{Langkah 3: Menghitung selisih:}
    \[
    g(1+h) - g(1) = (h^2 + 5h + 4) - 4 = h^2 + 5h = h(h + 5) \quad \text{\textbf{[4 poin]}}
    \]
    \item \textbf{Langkah 4: Membagi dengan $h$:}
    \[
    \frac{g(1+h) - g(1)}{h} = h + 5 \quad \text{\textbf{[4 poin]}}
    \]
    \item \textbf{Langkah 5: Evaluasi limit saat $h \to 0$:}
    \[
    g'(1) = \lim_{h \to 0}(h + 5) = 5 \quad \text{\textbf{[3 poin]}}
    \]
\end{itemize}

\paragraph{(b) [25 Poin] Turunan Implisit, Garis Singgung, dan Hampiran Linear}
Kurva: $x^2 + 2\sin(y) = 9$ di titik $(3, 0)$.
\begin{enumerate}[label=(\roman*), leftmargin=18pt]
    \item \textbf{[10 Poin] Turunan implisit $\dfrac{dy}{dx}$ di $(3, 0)$:}
    \begin{align*}
        \frac{d}{dx}\left[ x^2 + 2\sin(y) \right] &= \frac{d}{dx}[9] \\
        2x + 2\cos(y)\frac{dy}{dx} &= 0 \quad \text{\textbf{[5 poin]}} \\
        \frac{dy}{dx} &= \frac{-2x}{2\cos(y)} = \frac{-x}{\cos(y)} \quad \text{\textbf{[2 poin]}}
    \end{align*}
    Substitusikan titik $(3, 0)$:
    \[
    \left.\frac{dy}{dx}\right|_{(3, 0)} = \frac{-3}{\cos(0)} = \frac{-3}{1} = -3 \quad \text{\textbf{[3 poin]}}
    \]

    \item \textbf{[7 Poin] Persamaan garis singgung di $(3, 0)$:}
    Gradien $m = -3$.
    \begin{align*}
        y - y_0 &= m(x - x_0) \\
        y - 0 &= -3(x - 3) \implies y = -3x + 9 \quad \text{\textbf{[7 poin]}}
    \end{align*}

    \item \textbf{[8 Poin] Taksiran nilai $y$ ketika $x = 3{,}02$:}
    \begin{itemize}
        \item \textit{Melalui Hampiran Linear:}
        \[
        L(x) = -3x + 9 \implies y(3{,}02) \approx -3(3{,}02) + 9 = -9{,}06 + 9 = -0{,}06 \quad \text{\textbf{[8 poin]}}
        \]
        \item \textit{Melalui Diferensial:}
        $\Delta x = dx = 0{,}02$.
        \[
        dy = (-3)(0{,}02) = -0{,}06 \implies y(3{,}02) \approx 0 + (-0{,}06) = -0{,}06.
        \]
    \end{itemize}
\end{enumerate}

\vspace{10pt}
\subsection*{Soal 2 [Bobot: 55 Poin]}
Fungsi: $g(x) = x^3 - 12x$ pada selang tertutup $[-3, 3]$.

\paragraph{(a) [20 Poin] Kemonotonan \& Ekstrem Lokal}
\begin{itemize}[leftmargin=18pt]
    \item \textbf{Turunan pertama:}
    \[
    g'(x) = 3x^2 - 12 = 3(x^2 - 4) = 3(x-2)(x+2) \quad \text{\textbf{[5 poin]}}
    \]
    Titik stasioner: $x = -2$ dan $x = 2$.
    \item \textbf{Uji tanda $g'(x)$:}
    \begin{itemize}
        \item Pada $(-3, -2)$: $g'(x) > 0 \implies g$ \textbf{monoton naik} pada $[-3, -2]$.
        \item Pada $(-2, 2)$: $g'(x) < 0 \implies g$ \textbf{monoton turun} pada $[-2, 2]$.
        \item Pada $(2, 3)$: $g'(x) > 0 \implies g$ \textbf{monoton naik} pada $[2, 3]$.
    \end{itemize}
    (\textit{Menyatakan selang kemonotonan dengan benar: \textbf{[8 poin]}})
    \item \textbf{Ekstrem lokal:}
    \begin{itemize}
        \item Di $x = -2$: tanda $g'$ berganti dari $+$ ke $-$, maka ada \textbf{maksimum lokal} di $x = -2$ dengan nilai $g(-2) = (-2)^3 - 12(-2) = -8 + 24 = 16$. \textbf{[3.5 poin]}
        \item Di $x = 2$: tanda $g'$ berganti dari $-$ ke $+$, maka ada \textbf{minimum lokal} di $x = 2$ dengan nilai $g(2) = 2^3 - 12(2) = 8 - 24 = -16$. \textbf{[3.5 poin]}
    \end{itemize}
\end{itemize}

\paragraph{(b) [15 Poin] Kecekungan \& Titik Belok}
\begin{itemize}[leftmargin=18pt]
    \item \textbf{Turunan kedua:}
    \[
    g''(x) = \frac{d}{dx}[3x^2 - 12] = 6x \quad \text{\textbf{[4 poin]}}
    \]
    $g''(x) = 0 \iff x = 0$.
    \item \textbf{Selang Kecekungan:}
    \begin{itemize}
        \item Pada $(-3, 0)$: $g''(x) < 0 \implies$ kurva \textbf{cekung ke bawah}. \textbf{[3 poin]}
        \item Pada $(0, 3)$: $g''(x) > 0 \implies$ kurva \textbf{cekung ke atas}. \textbf{[3 poin]}
    \end{itemize}
    \item \textbf{Titik Belok:}
    Koordinat titik belok di $(0, g(0)) = (0, 0)$. \textbf{[5 poin]}
\end{itemize}

\paragraph{(c) [20 Poin] Titik Kritis, Nilai Maksimum Global, \& Nilai Minimum Global pada $[-3, 3]$}
\begin{itemize}[leftmargin=18pt]
    \item \textbf{Titik kritis pada $[-3, 3]$:}
    Titik ujung: $x = -3, 3$. Titik stasioner: $x = -2, 2$. Himpunan titik kritis: $\{-3, -2, 2, 3\}$. \textbf{[6 poin]}
    \item \textbf{Evaluasi nilai fungsi pada seluruh titik kritis:}
    \begin{align*}
        g(-3) &= (-3)^3 - 12(-3) = -27 + 36 = 9 \quad \text{\textbf{[2 poin]}} \\
        g(-2) &= 16 \quad \text{\textbf{[2 poin]}} \\
        g(2)  &= -16 \quad \text{\textbf{[2 poin]}} \\
        g(3)  &= 3^3 - 12(3) = 27 - 36 = -9 \quad \text{\textbf{[2 poin]}}
    \end{align*}
    \item \textbf{Kesimpulan Ekstrem Global:}
    \begin{itemize}
        \item \textbf{Nilai Maksimum Global} adalah \textbf{16} (tercapai di $x = -2$). \textbf{[3 poin]}
        \item \textbf{Nilai Minimum Global} adalah \textbf{-16} (tercapai di $x = 2$). \textbf{[3 poin]}
    \end{itemize}
\end{itemize}

\end{document}
